% Stage 8 section framing — 2.11 Summary
\section{Summary}
\label{sec:ch02.11}
% CH02-NARR-EXP:02.11
Maxwell evolution generates fields; outgoing selection determines export; finite instruments recover only part of that export \cite{maxwell1873,stratton1941,sommerfeld1949}. Vector spherical harmonics, source realizability, and interaction formalism beyond this chapter are cited in Chapters~3--4 without local re-derivation here \cite{harrington2001,devaney2004}.
Chapter~\ref{ch:02} began from a premise that Maxwell's equations, once supplemented by
outgoing radiation conditions, do not merely \emph{propagate} fields---they \emph{select} which
field content qualifies as radiative export. That selection is implemented by the radiation
operator $\mathcal{R}_\omega$, whose image is what finite instruments ultimately interrogate,
always incompletely and always through additional rank limits imposed by geometry and hardware.
The distinction is not pedantic. A near-field scanner that records thirty-six complex samples
has performed a serious measurement; yet the number answers a question about accessible angular
structure only after one traces which term in the master degree-of-freedom partition removed
the unresolved content---source nullity, aperture rank, evanescent storage, spectral compression,
or geometric suppression.


The four worked configurations in Section~\ref{sec:ch210} were composed to make that partition
tangible. The electrically small dipole showed that modal truncation error tracks compact singular
values, not non-radiating nullity. The loop-superposition problem showed that far patterns can
remain invariant while near probes register a thirty-five percent increment in tangential
amplitude. The evanescent-overlap problem showed that near-shell energy can exceed the
radiative component by more than a factor of two while far amplitude $a(\widehat{\mathbf{k}})$
remains fixed, and that applying far-field extractors to raw $\mathcal{S}_\omega[\Jimp]$ can
bias angular export by forty percent. The structured ring array showed that azimuthal phase
organization in $a_\sigma(\widehat{\mathbf{k}})$ can coexist with vanishing signed angular flux
until a modest amplitude imbalance restores transport. Each result is a laboratory caution stated
in operator language: existence of field does not imply export; export does not imply
wattmeter reading; $\oint_{\Gamma_R}$ closure does not imply unique source identification
\cite{maxwell1873}.
Read against that laboratory background, the chapter's central lesson is a completion criterion,
not a slogan. \ref{ch:01} established $\mathcal{D}_J=1$---what impressed currents and phasor
fields may enter the monograph at all. Chapter~\ref{ch:02} establishes $\mathcal{E}_{\mathrm{acc}}^{(N)}$: which
radiative content survives outgoing selection and which parts finite geometry and finite
instruments can resolve. Throughout, qualified reporting applies $\mathcal{R}_\omega$ before
Stokes parameters, Poynting vectors, or angular-flux functionals, and composes with regime
qualification inherited from \ref{ch:01} so that reactive shells are not mistaken for
transport \cite{harrington2001}. From the vector-network-analyzer reciprocity checks of Section~\ref{sec:ch28} through the directional
decompositions of Section~\ref{sec:ch29} to the four worked configurations of Section~\ref{sec:ch210},
the chapter supplies operator criteria for diagnosing pattern residuals by mechanism rather than by fitting alone.
